\documentclass[a4paper,UKenglish,cleveref, autoref, thm-restate]{lipics-v2021}

\usepackage{agda}
\usepackage{unicode}
\usepackage[utf8]{inputenc}
\usepackage{tikz-cd}

%TODO remove there two
\usepackage{silence}
\WarningFilter{latex}{Font shape}

\bibliographystyle{plainurl}% the mandatory bibstyle

\title{Dummy title} %TODO Please add

%\titlerunning{Dummy short title} %TODO optional, please use if title is longer than one line

\author{Jane {Open Access}}{Dummy University Computing Laboratory, [optional: Address], Country \and My second affiliation, Country \and \url{http://www.myhomepage.edu} }{johnqpublic@dummyuni.org}{https://orcid.org/0000-0002-1825-0097}{(Optional) author-specific funding acknowledgements}%TODO mandatory, please use full name; only 1 author per \author macro; first two parameters are mandatory, other parameters can be empty. Please provide at least the name of the affiliation and the country. The full address is optional. Use additional curly braces to indicate the correct name splitting when the last name consists of multiple name parts.

\author{Joan R. Public\footnote{Optional footnote, e.g. to mark corresponding author}}{Department of Informatics, Dummy College, [optional: Address], Country}{joanrpublic@dummycollege.org}{[orcid]}{[funding]}

\authorrunning{J. Open Access and J.\,R. Public} %TODO mandatory. First: Use abbreviated first/middle names. Second (only in severe cases): Use first author plus 'et al.'

\Copyright{Jane Open Access and Joan R. Public} %TODO mandatory, please use full first names. LIPIcs license is "CC-BY";  http://creativecommons.org/licenses/by/3.0/

\ccsdesc[100]{\textcolor{red}{Replace ccsdesc macro with valid one}} %TODO mandatory: Please choose ACM 2012 classifications from https://dl.acm.org/ccs/ccs_flat.cfm 

\keywords{containers, type theory, semantics} %TODO mandatory; please add comma-separated list of keywords

\category{} %optional, e.g. invited paper

\relatedversion{} %optional, e.g. full version hosted on arXiv, HAL, or other respository/website
%\relatedversiondetails[linktext={opt. text shown instead of the URL}, cite=DBLP:books/mk/GrayR93]{Classification (e.g. Full Version, Extended Version, Previous Version}{URL to related version} %linktext and cite are optional

%\supplement{}%optional, e.g. related research data, source code, ... hosted on a repository like zenodo, figshare, GitHub, ...
%\supplementdetails[linktext={opt. text shown instead of the URL}, cite=DBLP:books/mk/GrayR93, subcategory={Description, Subcategory}, swhid={Software Heritage Identifier}]{General Classification (e.g. Software, Dataset, Model, ...)}{URL to related version} %linktext, cite, and subcategory are optional

%\funding{(Optional) general funding statement \dots}%optional, to capture a funding statement, which applies to all authors. Please enter author specific funding statements as fifth argument of the \author macro.

\acknowledgements{I want to thank \dots}%optional

\nolinenumbers %uncomment to disable line numbering

%Editor-only macros:: begin (do not touch as author)%%%%%%%%%%%%%%%%%%%%%%%%%%%%%%%%%%
\EventEditors{John Q. Open and Joan R. Access}
\EventNoEds{2}
\EventLongTitle{42nd Conference on Very Important Topics (CVIT 2016)}
\EventShortTitle{CVIT 2016}
\EventAcronym{CVIT}
\EventYear{2016}
\EventDate{December 24--27, 2016}
\EventLocation{Little Whinging, United Kingdom}
\EventLogo{}
\SeriesVolume{42}
\ArticleNo{23}
%%%%%%%%%%%%%%%%%%%%%%%%%%%%%%%%%%%%%%%%%%%%%%%%%%%%%%

% \usepackage{minted}

\begin{document}

\maketitle

\begin{abstract}
This is abstract  
\end{abstract}

\section{Introduction}
\label{sec:introduction}

An introduction.

\section{Background and Metatheory}
\label{sec:metatheory}

Our metatheory is Martin-L\"of type theory. Propositional equality \AgdaDatatype{\_≡\_} is represented exclusively by the identity type, which is eliminated by the \emph{J}-rule. We do not assume uip - the uniqueness of identity proofs, also known as the \emph{K}-rule, and we impose neither higher inductive structure nor truncation on any of the types considered.

Under these assumptions, the categories studied in this paper are wild categories. Each consists of a class of objects and a \emph{class} of morphisms, equipped with identity morphisms and a composition operation that satisfy the left and right unit laws and associativity. Since identity types are not assumed to be propositions, the collections of morphisms are not required to be h-Sets.

It is natural to formalise our development in Agda with \AgdaPragma{--without-K} flag. Accordingly, we postulate extensional equality for function types - \AgdaFunction{funExt}, together with its computation rule - \AgdaFunction{funExt-β}.

\section{Basic Idea}
\label{sec:basic}

In this section, we introduce the denotational semantics of simple inductive types through a worked example, and then show how the construction generalises to containers and W-types. We then introduce
the category of containers.

Informally, the denotational semantics of inductive types is a correspondence between strictly-positive functors and their least fixed points, the latter being representable as strictly-positive inductive types. For example, the $\mathsf{Maybe}$ functor, defined as the operation $X \mapsto \mathbf{1} + X$ on \AgdaPrimitive{Type}, has the natural numbers \AgdaDatatype{ℕ} as its least fixed point.

Formally, given any \AgdaBound{A : }\AgdaPrimitive{Type}, and any F-algebra \AgdaBound{α :} \AgdaPrimitive{Maybe }\AgdaBound{A → A}, there is an unique folding function \AgdaFunction{foldℕ}\AgdaSpace{}%
\AgdaSymbol{:}\AgdaSpace{}%
\AgdaSymbol{(}\AgdaFunction{Maybe}\AgdaSpace{}%
\AgdaGeneralizable{A}\AgdaSpace{}%
\AgdaSymbol{→}\AgdaSpace{}%
\AgdaGeneralizable{A}\AgdaSymbol{)}\AgdaSpace{}%
\AgdaSymbol{→}\AgdaSpace{}%
\AgdaDatatype{ℕ}\AgdaSpace{}%
\AgdaSymbol{→}\AgdaSpace{}%
\AgdaGeneralizable{A}, such makes the following diagram commutes:

\[
\begin{tikzcd}[column sep=large, row sep=large]
  \mathsf{Maybe}\,\mathbb{N}
    \arrow[r, "{[\mathsf{zero},\,\mathsf{suc}]}"]
    \arrow[d, "\mathsf{Maybe}\,(\mathsf{fold}\,\alpha)"'] &
  \mathbb{N}
    \arrow[d, dashed, "\mathsf{fold}\,\alpha"] \\
  \mathsf{Maybe}\,A
    \arrow[r, "\alpha"'] &
  A
\end{tikzcd}
\]

To formalise this explicitly, it suffices to construct two functions: \AgdaFunction{commuteℕ}, which witnesses that \AgdaFunction{foldℕ} is a homomorphism of $\mathsf{Maybe}$-algebras, and \AgdaFunction{uniquefoldℕ}, which proves the uniqueness. Together, these establish that $\mathbb{N}$ is the initial $\mathsf{Maybe}$-algebra.

\paragraph*{Containers}
\begin{itemize}
  \item intuition, definition, extension
  \item 
\end{itemize}

\section{Second-Order}

This is second-order.

\section{Higher-Order}

This is higher-order.

\section{Conclusion and Future Work}

This is conclusion and future work.

%%
%% Bibliography
%%

\bibliography{references}

\end{document}
